% =====================================================================
%  CARNET D'ENTRAINEMENT — FICHE 1 : Les chiffres significatifs
%  THÈME 2 — Notation scientifique & changements d'unité (nombres exacts)
%  Fiche TUTO
% =====================================================================
\documentclass[11pt,a4paper]{article}
\usepackage[utf8]{inputenc}
\usepackage[T1]{fontenc}
\usepackage{lmodern}
\usepackage[french]{babel}
\usepackage{amsmath,amssymb,mathtools}
\usepackage{icomma}
\usepackage{siunitx}
\sisetup{output-decimal-marker={,},per-mode=symbol}
\usepackage[version=4]{mhchem}
\usepackage{xcolor}
\usepackage{tikz}
\usepackage[most]{tcolorbox}
\usepackage{enumitem}
\usepackage{booktabs}
\usepackage{array}
\usepackage{fancyhdr}
\usepackage[a4paper,margin=2cm,headheight=26pt,headsep=8pt]{geometry}
\usetikzlibrary{arrows.meta,calc}

\definecolor{primary}{RGB}{150,98,162}
\definecolor{primarydark}{RGB}{92,52,110}
\definecolor{formulacol}{RGB}{198,93,9}
\definecolor{reglecol}{RGB}{150,98,162}
\definecolor{methcol}{RGB}{90,90,100}
\definecolor{excol}{RGB}{0,135,90}
\definecolor{attncol}{RGB}{200,40,55}
\newcommand{\cs}{chiffre significatif}
\newcommand{\CS}{chiffres significatifs}

\pagestyle{fancy}\fancyhf{}
\fancyhead[L]{\small\textcolor{primarydark}{\textbf{Fiche 1} — Chiffres significatifs}}
\fancyhead[R]{\small\textcolor{primarydark}{\textsc{Tuto · Thème 2}}}
\fancyfoot[C]{\small\thepage}
\renewcommand{\headrulewidth}{0.8pt}
\renewcommand{\headrule}{\hbox to\headwidth{\color{primary}\leaders\hrule height \headrulewidth\hfill}}

\tcbset{boxbase/.style={enhanced,breakable,boxrule=0.9pt,arc=2.5pt,
  left=3mm,right=3mm,top=1.6mm,bottom=1.6mm,fonttitle=\bfseries,coltitle=white}}
\newtcolorbox{regle}[1][]{boxbase,colback=reglecol!7,colframe=reglecol,
  title={\textsc{Règle}\ifx\relax#1\relax\relax\else\ ---\ \textmd{\itshape #1}\fi}}
\newtcolorbox{formule}[1][]{boxbase,colback=formulacol!7,colframe=formulacol,
  title={\textsc{Définition}\ifx\relax#1\relax\relax\else\ ---\ \textmd{\itshape #1}\fi}}
\newtcolorbox{methode}[1][]{boxbase,colback=methcol!7,colframe=methcol,sharp corners,
  title={\textsc{Méthode}\ifx\relax#1\relax\relax\else\ : \textmd{\itshape #1}\fi}}
\newtcolorbox{exemple}[1][]{boxbase,colback=excol!6,colframe=excol,
  title={\textsc{Exemple}\ifx\relax#1\relax\relax\else\ : \textmd{\itshape #1}\fi}}
\newtcolorbox{attention}[1][]{boxbase,colback=attncol!6,colframe=attncol,
  title={\textsc{Attention}\ifx\relax#1\relax\relax\else\ : \textmd{\itshape #1}\fi}}

\setlength{\parindent}{0pt}\setlength{\parskip}{2pt}

\begin{document}

\begin{center}
{\LARGE\bfseries\color{primarydark} Thème 2 — Notation scientifique \& unités}\\[1pt]
{\large\color{primary} Écrire un nombre sans jamais changer sa précision}
\end{center}
\vspace{1mm}

La notation scientifique est l'écriture reine de la chimie : elle gère aussi bien les très grands
nombres ($N_A\approx\qty{6,022e23}{\per\mole}$) que les très petits ($[\ce{H+}]=\qty{1,0e-7}{\mol\per\litre}$),
\emph{tout en affichant sans ambiguïté le nombre de \CS}. Ce thème entraîne deux gestes jumeaux :
\textbf{mettre en notation scientifique} et \textbf{changer d'unité}, sans jamais modifier la précision.

\begin{formule}[notation scientifique]
\[ \boxed{\;a\times 10^{\,n}\;}\qquad\text{avec}\quad 1\le |a| < 10\ \text{et}\ n\in\mathbb{Z}. \]
$a$ est la \textbf{mantisse} : sa partie entière est \emph{un seul} chiffre non nul, et elle porte
\emph{tous} les \CS. $10^{\,n}$ fixe l'\textbf{ordre de grandeur} et n'apporte aucun \cs.
\end{formule}

\begin{methode}[mettre en notation scientifique]
\begin{enumerate}[leftmargin=1.6em,topsep=1pt,itemsep=0pt]
  \item placer la virgule \emph{juste après le premier chiffre non nul} ;
  \item compter de combien de rangs la virgule s'est déplacée : c'est $|n|$ ($n>0$ si le nombre était
        $\ge 10$, $n<0$ s'il était $<1$) ;
  \item ne garder dans la mantisse que les \CS\ (les zéros de cadrage disparaissent).
\end{enumerate}
\end{methode}

\begin{regle}[le déplacement de la virgule ne change jamais le nombre de \CS]
Changer d'unité (\si{\milli\litre}$\leftrightarrow$\si{\litre}, \si{\gram}$\leftrightarrow$\si{\kilo\gram}\dots)
ou passer en notation scientifique, c'est \emph{seulement déplacer la virgule}. Cela ne modifie
\textbf{jamais} la précision : le nombre de \CS\ reste identique. Seuls apparaissent ou disparaissent
des \emph{zéros de cadrage}, qui ne comptent pas.
\end{regle}

\begin{exemple}[un même volume, 4 \CS\ dans toutes les écritures]
\[ V=\underbrace{\num{20,96}}_{4\ \text{CS}}\,\si{\milli\litre}
    =\underbrace{\num{0,02096}}_{4\ \text{CS}}\,\si{\litre}
    =\underbrace{\num{2,096e-2}}_{4\ \text{CS}}\,\si{\litre}. \]
Les deux zéros de tête de $\num{0,02096}$ ne sont pas significatifs : ils disparaissent en notation
scientifique. La précision, elle, n'a pas bougé.
\end{exemple}

\begin{regle}[nombres exacts : une infinité de \CS]
Un nombre \textbf{exact} provient d'un \emph{comptage} ($2$ atomes de \ce{H} dans \ce{H2O}, $3$ fioles)
ou d'une \emph{définition} ($\qty{1}{\kilo\gram}=\qty{1000}{\gram}$, $\qty{1}{\hour}=\qty{3600}{\second}$,
$\qty{1}{\centi\metre\cubed}=\qty{1}{\milli\litre}$). On lui attribue une \textbf{infinité} de \CS : il
ne limite \emph{jamais} la précision d'un résultat. Les constantes ($\pi$, $\sqrt2$) sont prises avec
\emph{au moins} autant de \CS\ que la donnée la plus précise.
\end{regle}

\begin{attention}[erreurs classiques de mantisse]
$\num{12,5e3}$ et $\num{0,4e5}$ ne sont \emph{pas} en notation scientifique : la mantisse doit vérifier
$1\le|a|<10$. On corrige : $\num{12,5e3}=\num{1,25e4}$ et $\num{0,4e5}=\num{4e4}$.
\end{attention}

\end{document}
